\chapter{Les fondements de la géométrie plane}
\label{exercices:12}
\section*{Exercice 12.1 -- Angles d'un triangle}
Un triangle $ABC$ est isocèle en $A$ et $\widehat A=44\degre$. Calculer les autres angles. Un triangle a deux angles de $35\degre$ et $55\degre$ : quelle est sa nature ? Justifier.

\section*{Exercice 12.2 -- Construire et décider si c'est possible}
Décrire la construction d'un triangle de côtés $5$, $7$ et $10\unite{cm}$. Pourquoi les longueurs $3$, $4$ et $8\unite{cm}$ ne permettent-elles pas un triangle ?

\section*{Exercice 12.3 -- Médiatrice et cercle}
Construire la médiatrice de $[AB]$ avec règle et compas. Un point $M$ vérifie $MA=MB$ et $MB=MC$. Que peut-on conclure si $A,B,C$ ne sont pas alignés ?

\section*{Exercice 12.4 -- Quel quadrilatère ?}
Les diagonales d'un quadrilatère non aplati ont le même milieu. Elles mesurent toutes deux $10\unite{cm}$. Quelle nature peut-on établir ? Si l'on apprend en plus qu'elles sont perpendiculaires, que devient la conclusion ?
